\Needspace{6\baselineskip}
\begin{center}\hypertarget{native-265-article191}{}\pdfbookmark[2]{191.}{article-191}191.\end{center}
Die Entwicklung der Theorie der Störungen aus gegebenen Elementen liegt
außerhalb unserer Aufgabe. Hier ist nur zu zeigen, wie sich eine
Näherungsbahn so korrigieren lässt, dass sie zusammen mit den gegebenen
Störungen den Beobachtungen möglichst genau entspricht. Am einfachsten
lässt sich diese Aufgabe mit einem Verfahren lösen, das den in
Art.~124, 165, 188 dargelegten Verfahren ähnelt. Für die Zeitpunkte
aller Beobachtungen, die man zu diesem Zweck verwenden will und deren
Zahl je nach den Umständen drei, vier oder mehr betragen kann, berechne
man die Zahlenwerte der Störungen aus den zugehörigen Gleichungen, und
zwar für die Längen auf der Bahn, die Radiusvektoren und die heliozentrischen Breiten.
Die Argumente für diese Rechnung entnimmt man den elliptischen
Näherungselementen, auf denen die Theorie der Störungen beruhte. Dann
wähle man aus allen Beobachtungen zwei aus und berechne für sie die
Entfernungen von der Erde aus denselben Näherungselementen. Diese
Entfernungen bilden die erste Hypothese. Die zweite und die dritte
Hypothese erhält man durch geringfügige Änderung dieser Entfernungen.
Für jede Hypothese bestimme man dann aus den beiden geozentrischen
Örtern die heliozentrischen Positionen und die Entfernungen von der
Sonne. Aus den heliozentrischen Positionen leitet man, nachdem die
Breiten um die Störungen bereinigt worden sind, die Länge des
aufsteigenden Knotens, die Bahnneigung und die Längen auf der Bahn ab.
Bei dieser Rechnung muss das Verfahren aus Art.~110 etwas abgeändert
werden, wenn es sich zu lohnen scheint, die säkulare Änderung der
Knotenlänge und der Neigung zu berücksichtigen. Es seien $\beta,\beta'$
die um die periodischen Störungen bereinigten heliozentrischen Breiten,
$\lambda,\lambda'$ die heliozentrischen Längen, $\AscendingNode,\AscendingNode+\Delta$
die Längen des aufsteigenden Knotens und $i,i+\delta$ die Bahnneigungen.
Die Gleichungen schreibt man dann zweckmäßig in folgender Form.
\[
\begin{gathered}
\mathrm{tang}\,\beta=\mathrm{tang}\,i\sin(\lambda-\AscendingNode)\\
\frac{\mathrm{tang}\,i}{\mathpunct{.}\mathrm{tang}(i+\delta)}\mathrm{tang}\,\beta'
=\mathrm{tang}\,i\sin(\lambda'-\Delta-\AscendingNode).
\end{gathered}
\]
\SourceLeaf{266}Den Wert von $\frac{\mathrm{tang}\,i}{\mathrm{tang}(i+\delta)}$
erhält man mit aller nötigen Genauigkeit, indem man für $i$ einen
Näherungswert einsetzt. Dann lassen sich $i$ und $\AscendingNode$ nach den
üblichen Verfahren ermitteln.

Von den beiden Längen auf der Bahn und den beiden Radiusvektoren zieht man
ferner die Summen der Störungen ab, um die rein elliptischen Werte zu
erhalten. Dabei ist auch die Wirkung der säkularen Änderungen der
Lage des Perihels und der Exzentrizität auf die Länge auf der Bahn und
den Radiusvektor zu berücksichtigen. Sie ist anhand der
Differentialformeln aus Abschnitt~I des ersten Buches zu bestimmen und
gleich mit den periodischen Störungen zu verbinden, wenn die
Beobachtungen so weit auseinanderliegen, dass sich die Berücksichtigung
dieser Wirkung lohnt. Aus den so korrigierten Längen auf der Bahn und
Radiusvektoren sowie den entsprechenden Zeitpunkten bestimmt man die
übrigen Elemente. Schließlich berechnet man aus diesen Elementen die
geozentrischen Positionen für alle übrigen Beobachtungen. Durch den
Vergleich mit den beobachteten Positionen erhält man auf die in
Art.~188 erläuterte Weise das System von Entfernungen, aus dem sich
Elemente ergeben, die allen übrigen Beobachtungen möglichst gut
entsprechen.
